\section{RRAM：原生共享读线与完整二态更新}

本例采用TaO$_x$ WH-2T1R二状态参考路径，八个等尺度bit平面组成 $W[64,128]$ INT8矩阵，有效容量8192 Byte。读侧保留32项活动及原生子阵列选通；写侧配置128个实际驱动、256个窗口比较器，完整RESET、masked SET、读验与偏置退出均计时。典型 $\rho\approx25.0$ MB/s、$\tau\approx4.29$ MB/s、$\RI^*\approx5.82$，对应一次完整矩阵装载的持续服务。

\subsection{原生组织、编码与读服务}

\sid{RRAM-05} PDF p.3的Fig.3给出64行128列宏，四个64行32列子阵列共享按输出行布置的TBL，CIMSEL选择活动子阵列；memory模式有独立BL/SL选通，CIM模式走T2及TBL。Fig.4中的IO器件T1承受写入，core器件T2处理计算。本参考采用八个完整的64行128列局部宏，每个宏承载一个权重bit平面，全部使用 $m=1$ 等尺度T2。四个32输入子阵列依次求值，避免把共享TBL免费拆成独立位平面。该组合共有65536个数据cell，不用互补编码或原文 $m=1/2/4$ 空间加权来改变逻辑payload。

64个输出按16个一组读出；每平面16个ADC，总计128个，16条数字输出通道完成位权合并及累加。输入寄存器1024 bit，64个输出使用23-bit容器。32项二态部分和有33个理想等级，仍采用公共约8有效bit转换及校准近似合同；权重和输入最高位按负位权重构。原文5 ns的PH0及本地判决量级用于前端建立锚点，不把板级完整四位66 ns周期重复叠加到独立SAR阶段。

每个128 Byte输入向量需要八个输入bit、四个输入子阵列、四个输出组，共128轮求值和转换，256个数字拍。因此
\begin{equation}
\Delta_S=128(T_I+5\ns+T_A+2T_D)+2T_D.
\label{05_rram:eq:read}
\end{equation}
典型间隔为5.130 $\mu$s。既有位平面、输入重复施加和数字部分和不增加逻辑输入字节；输出完成包含在服务中。

\subsection{128个实际写目标与完整装载覆盖}

写入沿原memory BL/SL路径进行。每平面在同一输出WL、同一32列子阵列中选16列：两个16列组各选8列。每列有独立BL/SL驱动，整批施加同一RESET极性或按目标mask执行SET；未选列禁止写入，TBL计算通路隔离，不存在同列同时要求相反电压的冲突。八平面因此同时更新128个物理目标，恰好完成16个INT8权重。下半事务选其余16列；64个输出、四个输入组和两个半组构成512个事务，逐一覆盖8192个权重，无遗漏或重复。

该配置采用128个耐压/限流驱动和256个窗口比较器，不声称原实测宏已具有这些资源。每lane额定300 $\mu$A，更新域总额定38.4 mA；支持最高1.8 V编程偏置，需要专用IO、电平转换及供给。原文core/IO测量供电不能替代此耐压条件。128-bit目标缓冲、128-bit done/mask、双组地址控制与独立比较器在完整周期内占用；比较判决锁存在sense时隙内，另付一拍done归约与mask提交。16驱动、32比较器的缩资源配置单列，不作为RRAM材料并行上限。

LRS参考验收窗为8--12 k$\Omega$，HRS至少70 k$\Omega$，分别围绕\sid{RRAM-05}的10/100 k$\Omega$典型值选择。它们是二态窗口，不是原文保证界。每个物理目标先RESET，再对需要LRS的bit执行masked SET；禁用lane不耗编程电流，但本预算保留SET时隙，不靠数据稀疏跳过来增加吞吐。每次尝试都重新读验，未达窗lane保留在重试mask中；最终未达窗不发布完整矩阵。

窗口的计算作用来自WH-2T1R的转移通路：RRAM与T1分压决定T2栅压，亚阈值T2再产生TBL求和电流。普通memory能够区分HRS/LRS，尚不能单独保证32项MAC所需的ON电流与低漏电。参考窗用于约束可校准的二态工作域，并不存储8-bit模拟电导；在既定偏置下保持约0.5 $\mu$A的LRS基准和受控HRS贡献，是近似部分和服务的条件。Fig.4的典型电阻及Fig.7的变化模型没有给出8/12/70 k$\Omega$窗边界到电流的完整转移数据，因此不把这些阈值称为已验证的计算精度充分界。若SET后电阻低于8 k$\Omega$，继续同极性SET不能保证回窗；反极性救援及后续读验属于额外操作，不能按当前有限次数预算记为已成功装载。

\subsection{脉冲、局部切换与写rail生命周期}

\sid{RRAM-01} PDF p.10报告实际1 $\mu$s SET/RESET脉冲，以及受外部DAC/ADC控制限制的1--10 $\mu$s读回。其HfO$_x$/TaO$_x$堆栈不同于本参考TaO$_x$，编程目标上限为30/40 $\mu$S，验收容差为$\pm1\ \mu$S；本LRS窗则对应83--125 $\mu$S，超出该目标区间。因而1 $\mu$s保留为有条件的跨实现波形时隙，不保证当前窗口可达，也不由窗口较宽推导所需次数较少。多级8.52次平均脉冲、56 $\mu$s投影及成功率均不迁移。

\sid{RRAM-03} PDF p.5 E节及Fig.9报告1 Mb二态阵列多数cell按名义条件完成、数十个cell需要追加编程，但未报告绝对脉宽、追加次数分布或本例窗口。\sid{RRAM-06} PDF p.1及p.2 Figs.1、3--7给出RESET通常较慢、组内all-done、mask与超时机制，其归一化时间不能确定2/1。参考选点在两相一次完成的1/1基础上，为RESET多保留一次修正，取2/1；4/2表示更长的有限完成条件。三者均不是均值、中位数、良率分位或材料最坏上界。每批按最慢活动lane完成计时，已done lane可以关断，但不以提前完成提高表中吞吐。

写电源保持与局部偏置切换分开。整矩阵开始与结束各预留1 $\mu$s供给/rail建立和退出；512个事务期间rail保持，不在每次尝试重新上电。每个pulse之前仍要选择地址并建立目标BL/SL偏置，之后仍需隔离、退偏至新的verify条件；这两段局部切换分别记为 $g$，三情景取50/100/250 ns。

这些局部时间是固定驱动与负载条件下的工程预算。每lane的开关等效负载限制为1 pF；1.8 V摆幅需要1.8 pC。用300 $\mu$A额定减去 $1.8\mathrm{V}/8\mathrm{k}\Omega=225\ \mu$A的端态静态电流量级，剩余75 $\mu$A对应约24 ns充电量级。最短50 ns提供两倍这一量级，100/250 ns允许额外选通与建立余量。这个检查约束驱动、布线和供给，不把端态欧姆比值当作瞬态编程电流测量，也不证明实际开关可在该时间内完成。1 pF及相应时序是所选参考实现的条件，不是介质固有常数。

求值前端记为 $F_C=5$ ns，沿core T2、TBL及CIM感测通路；写后读验前端记为 $F_V=5$ ns，沿T1、IO、BL/SL及独立memory窗口比较器。两者具有不同通路、负载和判决终点，数值相等是参考选择，不能用PH0测量证明memory前端恰好也需5 ns。代码分别绑定两个参数；修改 $F_C$只改变求值，修改 $F_V$只改变读验。公共输入选通 $T_I$、数字拍 $T_D$同时用于两路；$T_B=T_A$是共同采样、判决和锁存能力的预算政策，故公共 $T_A$改变时两路同步重算，这一依赖不表示两路使用同一个ADC。

本地二态读验用 $T_I+F_V+T_B$，其中 $T_B=10/20/50$ ns覆盖窗口比较与锁存。随后一拍更新done/mask；pulse前地址控制另占一拍。外部1--10 $\mu$s测试读回不移作本地切换，低压感测也不取零。每完整矩阵的服务为
\begin{equation}
\begin{aligned}
t_{16}&=2T_D+(K_R+K_S)
 [1\us+2g+T_I+F_V+T_B+2T_D],\\
T_R&=512t_{16}+2\us,\qquad B_R=8192\ \mathrm{Byte}.
\end{aligned}
\label{05_rram:eq:write}
\end{equation}
128-bit首数据随命令捕获，准备和完成共两拍。$K_R/K_S$取1/1、2/1、4/2，描述整批最慢活动cell在该次数内达到窗口的正常完成情景；超出预算不伪装成成功。典型 $t_{16}=3.730\ \mu$s、$T_R=1.911760$ ms。独立16 Byte请求还需两端rail费用，为5.730 $\mu$s；它与整矩阵中嵌入事务不是同一延迟。

\begin{table}[htbp]\centering\small
\begin{tabular}{@{}lrrrrr@{}}\toprule
情景 & 脉冲 & 局部切换 & 新读验 & 控制/rail & 总计($\mu$s)\\\midrule
乐观 & 1.02e+03 & 102 & 17.4 & 8.14 & 1.15e+03\\
典型 & 1.54e+03 & 307 & 46.1 & 22.5 & 1.91e+03\\
悲观 & 3.07e+03 & 1.54e+03 & 200 & 73.7 & 4.88e+03\\
\bottomrule\end{tabular}
\caption{全矩阵更新占用分解，时间均为微秒。rail仅在整矩阵边界起退，局部切换每次尝试都保留。}\end{table}


典型写占用约80\%来自完整编程脉冲，约16\%来自局部偏置切换。rail边界费用只占一小部分，但在单个小请求中不能忽略。没有把RESET放在不计时的idle，也没有建立免费全矩阵shadow缓冲。

\subsection{三情景、资源对照与两表接口}

\begin{table}[htbp]\centering\small
\begin{tabular}{@{}lrrrrrr@{}}\toprule
情景 & $K_R/K_S$ & $\Delta_S$($\mu$s) & $T_R$(ms) & $\rho$(MB/s) & $\tau$(MB/s) & $\RI^*$\\\midrule
乐观 & 1/1 & 2.692 & 1.1520 & 47.5 & 7.11 & 6.69\\
典型 & 2/1 & 5.130 & 1.9118 & 25 & 4.29 & 5.82\\
悲观 & 4/2 & 10.900 & 4.8814 & 11.7 & 1.68 & 7\\
\bottomrule\end{tabular}
\caption{原生 $K=128,N=64$，完整8192 Byte装载含一次写rail起退；固定128驱动、256比较器。三情景为声明次数内完成的条件选点，不是统计快慢界。}\end{table}


主结果始终按完整矩阵计算 $\rho=128/\Delta_S$、$\tau=8192/T_R$，并得到 $U^*=T_R/\Delta_S=64\RI^*$。因此 $\tau$与 $U^*$使用相同rail及装载边界。位平面复用、CIMSEL分时和实际写lane决定两路服务，不能把16 Byte事务的稳态比例直接当任意小更新能力。

\begin{table}[htbp]\centering\small
\begin{tabular}{@{}lrrrr@{}}\toprule
资源 & 额定电流(mA) & $T_R$(ms) & $\tau$(MB/s) & $\RI^*$\\\midrule
W128 & 38.4 & 1.9118 & 4.29 & 5.82\\
W16 & 4.8 & 15.2442 & 0.537 & 46.4\\
\bottomrule\end{tabular}
\caption{W16缩资源对照将驱动与窗口比较器同比减少；同一脉冲/局部切换及完整rail生命周期，读侧不变。}\end{table}

\begin{table}[htbp]\centering\small
\begin{tabular}{@{}lrrr@{}}\toprule
$K_R/K_S$ & $T_R$(ms) & $\tau$(MB/s) & $\RI^*$\\\midrule
1/1 & 1.2769 & 6.42 & 3.89\\
2/1 & 1.9118 & 4.29 & 5.82\\
4/2 & 3.8164 & 2.15 & 11.6\\
\bottomrule\end{tabular}
\caption{固定参考外围与100 ns局部切换，仅改变尝试次数；不解释为概率分布或材料最坏值。}\end{table}


另将参考尝试次数及外围固定，单独扫描每段局部切换 $g=50/100/1000$ ns；机器数据保留完整结果，用于区分波形与外围假设。三情景保持128个驱动、38.4 mA额定及相同负载上限；没有将资源扩展混入快慢端。瞬时能力不扣除耐久次数，长期二态窗保持与使用环境另作可行性条件。

该点的位置主要由保留的原生共享TBL/32项分组，以及已明确付资源的128路写验和完整脉冲次数决定。跨stack脉冲达到本二态窗、同偏置电流校准及组内完成次数仍是参考实现的条件；数值范围不冒充同芯片配对测量或精确最坏界。独立扰动检查分别计数128个求值组、512个更新事务及典型每事务三次尝试，核对各前端、公共时隙、局部切换、脉冲和rail的传播斜率，并保持完整装载与 $U^*$相同边界。
